\chapter{Kéo Nói Vào Bài}
\chapterintro{Đừng để câu hỏi chỉ tạo ồn}{Một câu hỏi tốt phải kéo được lời nói quay vào nội dung.}

\section{Trong bài này, chỗ nào dễ hiểu nhất trước?}
\begin{cauhoibox}{Câu hỏi}
Trong bài này, chỗ nào dễ hiểu nhất trước?
\end{cauhoibox}
\begin{taisaocoich}
Hỏi cái dễ trước giúp học sinh có cảm giác mình làm được. Sau đó mới dễ kéo sang phần khó hơn.
\end{taisaocoich}
\begin{goiybien}
Rất hợp khi lớp đang sợ bài mới hoặc vừa học xong một đoạn dài.
\end{goiybien}
\ngancach

\section{Nếu phải giải thích cho bạn học chậm hơn em, em sẽ nói câu nào trước?}
\begin{cauhoibox}{Câu hỏi}
Nếu phải giải thích cho bạn học chậm hơn em, em sẽ nói câu nào trước?
\end{cauhoibox}
\begin{taisaocoich}
Câu này kéo học sinh từ chỗ hiểu cho mình sang chỗ diễn đạt cho người khác. Đó là mức hiểu sâu hơn.
\end{taisaocoich}
\begin{goiybien}
Có thể đổi thành: em sẽ vẽ gì trước, ví dụ gì trước, hay tránh nhầm chỗ nào trước.
\end{goiybien}
\ngancach

\section{Nếu bài này có một lỗi học sinh hay mắc, em đoán là lỗi gì?}
\begin{cauhoibox}{Câu hỏi}
Nếu bài này có một lỗi học sinh hay mắc, em đoán là lỗi gì?
\end{cauhoibox}
\begin{taisaocoich}
Học sinh thích đoán lỗi vì nó gần với trải nghiệm thật. Câu này giúp lớp chú ý hơn tới chỗ dễ nhầm mà không cần thầy cô giảng dài dòng trước.
\end{taisaocoich}
\begin{goiybien}
Sau 2 đến 3 câu trả lời, chốt bằng: \textit{``Đúng, hôm nay mình học để bớt nhầm ở chỗ đó.''}
\end{goiybien}
\ngancach

\section{Ý nào trong bài này chạm đời sống em nhất?}
\begin{cauhoibox}{Câu hỏi}
Ý nào trong bài này chạm đời sống em nhất?
\end{cauhoibox}
\begin{taisaocoich}
Câu hỏi kéo nội dung ra khỏi sách một chút rồi đưa ngược vào trải nghiệm của học sinh. Khi thấy bài có chỗ chạm, các em nói thật hơn.
\end{taisaocoich}
\begin{goiybien}
Hợp với văn, sử, giáo dục công dân, hoặc bất kỳ phần nào có thể nối vào đời sống gần của học sinh.
\end{goiybien}
\ngancach

\section{Nếu phải chọn ví dụ gần lớp mình nhất, em chọn ví dụ nào?}
\begin{cauhoibox}{Câu hỏi}
Nếu phải chọn ví dụ gần lớp mình nhất để hiểu bài này, em chọn ví dụ nào?
\end{cauhoibox}
\begin{taisaocoich}
Học sinh thường nói tốt hơn khi được lấy ví dụ từ thứ quanh mình. Câu hỏi này kéo lời nói vào bài mà không bị xa lạ.
\end{taisaocoich}
\begin{goiybien}
Cho 2 đến 3 ví dụ ngắn rồi chốt ví dụ gần nhất lên bảng để dùng trong cả tiết.
\end{goiybien}
